\documentclass[11pt]{article}

% ============================================================
% Packages
% ============================================================

\usepackage[margin=1in]{geometry}
\usepackage{amsmath,amssymb,mathtools}
\usepackage{enumitem}
\usepackage{hyperref}
\usepackage{xcolor}
\usepackage{microtype}

% ============================================================
% Macros
% ============================================================

\newcommand{\Pbb}{\mathbb{P}}
\newcommand{\Qbb}{\mathbb{Q}}
\newcommand{\E}{\mathbb{E}}
\newcommand{\R}{\mathbb{R}}
\newcommand{\Nbb}{\mathcal{N}}
\newcommand{\dd}{\,\mathrm{d}}
\newcommand{\1}{\mathbf{1}}

% Brownian / stochastic calculus
\newcommand{\F}{\mathcal{F}}
\newcommand{\Filtration}{(\mathcal{F}_t)_{t\ge 0}}

% ============================================================
% Environments
% ============================================================

\newtheorem{theorem}{Theorem}
\newtheorem{proposition}{Proposition}
\newtheorem{lemma}{Lemma}

\newtheorem{definition}{Definition}
\newtheorem{remark}{Remark}

% ============================================================
% Title
% ============================================================

\title{\textbf{Pricing and fundamental Principles}\\
\large Lecture Notes}
\author{Nathan Sauldubois}
\date{MSFE6233 --- Option Pricing}

\begin{document}
\maketitle
\vspace{-1.2em}


% ============================================================
\section{No-arbitrage principles and the Fundamental Theorems of Asset Pricing}
% ============================================================

\subsection{The dominance principle (first no-arbitrage principle)}

A very simple but fundamental no-arbitrage idea is the following:
if one contract always pays at least as much as another one, then it cannot cost less.

\paragraph{Set-up.}
Consider two contracts maturing at the same date \(T\), with payoffs \(A\) and \(B\),
and current prices \(a\) and \(b\), respectively.
If
$$
A \ge B
\qquad \text{in every state of the world,}
$$
then no-arbitrage requires
$$
\boxed{\ a \ge b.\ }
$$

\paragraph{Why?}
If instead \(a<b\), then one can:
\begin{itemize}[leftmargin=1.4em]
    \item buy the cheaper contract with larger payoff \(A\),
    \item sell the more expensive contract with smaller payoff \(B\).
\end{itemize}
The initial cashflow is
$$
-b+a<0,
$$
so the strategy generates a positive amount of cash at time \(0\),
while the terminal payoff is
$$
A-B\ge 0
$$
in every state.
This is an arbitrage. The slides present this as the first and most basic no-arbitrage argument. 
\paragraph{Interpretation.}
The dominance principle is the most elementary form of monotonicity of prices:
$$
\text{larger payoff} \quad \Longrightarrow \quad \text{larger price.}
$$
It is often the quickest way to derive pricing inequalities and bounds.

\subsection{Law of one price}

A second basic no-arbitrage principle is the \emph{law of one price}:

\begin{quote}
If two trading strategies generate exactly the same terminal cashflows in all states,
then they must have the same initial cost.
\end{quote}

Indeed, if two portfolios \(X\) and \(Y\) satisfy
$$
X_T=Y_T \quad \text{a.s.}
$$
but have different initial costs \(X_0\neq Y_0\), then buying the cheaper one and selling the more expensive one yields:
\begin{itemize}[leftmargin=1.4em]
    \item a strictly positive gain at time \(0\),
    \item zero payoff at maturity.
\end{itemize}
This is a pure arbitrage.

\paragraph{Why it matters for derivatives.}
If a derivative payoff \(H\) can be replicated by a self-financing strategy,
then its arbitrage-free price is uniquely determined:
it must equal the initial cost of the replicating portfolio.
This is the conceptual bridge between \emph{replication pricing} and \emph{martingale pricing}.

\subsection{Example: forward price from no-arbitrage}

\paragraph{Forward payoff.}
A forward contract delivering one share of stock at time \(T\) for a fixed delivery price \(K_t\)
has payoff
$$
S_T-K_t
$$
at maturity.

\paragraph{Replicating strategy (no dividends).}
The same payoff can be produced by:
\begin{itemize}[leftmargin=1.4em]
    \item buying one share today (cost \(S_t\)),
    \item borrowing the present value of \(K_t\), i.e.\ selling \(K_t\) units of the zero-coupon bond maturing at \(T\).
\end{itemize}
If \(P_t(T)\) denotes the price at time \(t\) of a zero-coupon bond paying \(1\) at time \(T\), then the initial cost is
$$
S_t-K_t P_t(T).
$$

Since a forward contract is entered at zero initial cost in equilibrium, no-arbitrage implies
$$
S_t-K_tP_t(T)=0,
$$
hence
$$
\boxed{\ K_t=\frac{S_t}{P_t(T)}. \ }
$$

\paragraph{Constant-rate case.}
If the short rate is constant \(r\), then \(P_t(T)=e^{-r(T-t)}\), and therefore
$$
\boxed{\ K_t=S_t e^{r(T-t)}. \ }
$$

\subsection{Arbitrage-free market and equivalent martingale measures}

The martingale pricing formula from Part A is not just a computational trick:
it is tied to the deep structure of arbitrage-free markets.

\paragraph{Informal principle.}
A market is arbitrage-free precisely when there exists a probability measure under which discounted traded asset prices are martingales.
Such a measure is called an \emph{equivalent martingale measure} (EMM).

This is exactly the role played by the risk-neutral measure \(\Qbb\) in the Black--Scholes model:
under \(\Qbb\), the discounted stock price \(D_tS_t\) is a martingale, and prices are computed by conditional expectation.

\paragraph{Why ``equivalent''?}
Equivalence means
$$
\Qbb \sim \Pbb,
$$
so \(\Pbb\) and \(\Qbb\) have the same null events.
In particular, the change of measure does not create or remove possible states of the world;
it only changes their probabilities.

\subsection{First Fundamental Theorem of Asset Pricing (informal statement)}

A standard form of the first fundamental theorem is:

\begin{theorem}[FTAP I --- informal]
A financial market is arbitrage-free if and only if there exists
an equivalent martingale measure \(\Qbb\) such that every discounted traded asset price process is a \(\Qbb\)-martingale
(or, in more general formulations, a local martingale).
\end{theorem}

\paragraph{How to read this theorem.}
\begin{itemize}[leftmargin=1.4em]
    \item \textbf{No arbitrage} is an economic condition.
    \item \textbf{Existence of an EMM} is a probabilistic condition.
    \item The theorem says these are two faces of the same idea.
\end{itemize}

\paragraph{Connection with your course.}
In Black--Scholes, Girsanov is precisely the tool used to construct this EMM:
starting from the physical measure \(\Pbb\), one changes probability so that the drift of the discounted stock disappears.

\subsection{Second Fundamental Theorem of Asset Pricing (informal statement)}

A second fundamental theorem links \emph{uniqueness} of the martingale measure to \emph{completeness} of the market.

\begin{theorem}[FTAP II --- informal]
In an arbitrage-free market, the following are essentially equivalent:
\begin{itemize}[leftmargin=1.4em]
    \item the market is complete (every claim can be replicated),
    \item the equivalent martingale measure is unique.
\end{itemize}
\end{theorem}

\paragraph{Interpretation.}
\begin{itemize}[leftmargin=1.4em]
    \item If the EMM is unique, then the arbitrage-free price of any replicable claim is uniquely determined.
    \item If there are many EMMs, then no-arbitrage alone typically gives only a range of admissible prices.
\end{itemize}

This matches the conclusion already emphasized in your notes:
$$
\text{complete market} \quad \Longrightarrow \quad \text{unique price by replication}.
$$

\subsection{Dominance, superhedging, and price bounds}

The dominance principle also leads naturally to \emph{superhedging} bounds.

\paragraph{Superhedging idea.}
Suppose a self-financing strategy with wealth \(X\) satisfies
$$
X_T \ge H
\qquad \text{a.s.}
$$
for a claim payoff \(H\).
Then \(X\) dominates the claim at maturity, so no-arbitrage implies that the initial cost of this strategy must be at least as large as the claim price:
$$
\boxed{\ V_0 \le X_0. \ }
$$

Similarly, if another strategy satisfies
$$
Y_T \le H,
$$
then its initial cost provides a lower bound:
$$
\boxed{\ Y_0 \le V_0. \ }
$$

\paragraph{Consequence.}
If one can find:
\begin{itemize}[leftmargin=1.4em]
    \item an upper replicating / superhedging strategy,
    \item and a lower subhedging strategy,
\end{itemize}
then one obtains rigorous no-arbitrage price bounds for the claim.

\paragraph{Special case: exact replication.}
If
$$
X_T=H,
$$
then both upper and lower bounds coincide, and the price is uniquely pinned down:
$$
V_0=X_0.
$$

\subsection{Practical takeaway for students}

The logical structure is:

\begin{enumerate}[leftmargin=1.6em]
    \item \textbf{No-arbitrage principles} (dominance, law of one price) rule out inconsistent prices.
    \item \textbf{Replication} gives exact prices when a payoff can be synthesized.
    \item \textbf{Equivalent martingale measures} translate no-arbitrage into a probabilistic statement.
    \item \textbf{Completeness} determines whether the price is unique.
\end{enumerate}

\paragraph{Big picture.}
The dominance principle is the first building block.
The fundamental theorems of asset pricing are the abstract global version of the same idea:
$$
\text{no-arbitrage}
\quad \Longleftrightarrow \quad
\text{martingale pricing under a suitable measure.}
$$

\begin{remark}
Pedagogically, the dominance principle is the simplest local no-arbitrage check,
whereas the fundamental theorems explain the full mathematical architecture behind risk-neutral valuation.
\end{remark}


\end{document}